\section*{Bab \ref{ch:b2:amines} --- Amina dan Senyawa Nitrogen}

\begin{solution}{exo:b2:amines:1}
Primer; sekunder; tersier; garam amonium kuaterner; primer (\omterm{def:b2:huckel:aromatic}{aromatik}).
\end{solution}

\begin{solution}{exo:b2:amines:2}
Etanamida $<$ anilina $<$ amonia $<$ metilamina ($\mathrm pK_a$ asam konjugasinya
$-0.6$, 4.60, 9.26, 10.66).
\end{solution}

\begin{solution}{exo:b2:amines:3}
\textit{N}-Metilsikloheksanimina, \ce{C6H10=N-CH3}: amina dan keton dalam
medium yang asam lemah (pH 4--5), dengan pengeluaran air yang terbentuk (zat
pengering atau distilasi azeotropik).
\end{solution}

\begin{solution}{exo:b2:amines:4}
Bromobenzena; benzonitril; fenol; 4-hidroksiazobenzena \ce{C6H5-N=N-C6H4-OH}
(kopling pada posisi para terhadap gugus \ce{O-} fenoksida).
\end{solution}

\begin{solution}{exo:b2:amines:5}
$\mathrm pK_a = 9.80$: pada pH 7 rasionya $[\ce{(CH3)3NH+}]/[\ce{(CH3)3N}] = 10^{2.8} \approx 630$:
bentuk terprotonasi dominan; air jeruk nipis (pH sekitar 2) memprotonasi
praktis seluruhnya. Mengocok larutan organik sebuah amina dengan asam berair
encer memindahkan amina itu, sebagai ion amoniumnya, ke dalam air; basa
membebaskannya kembali.
\end{solution}

\begin{solution}{exo:b2:amines:6}
1-(Sikloheks-1-en-1-il)pirolidina. Setelah adisi dan lepasnya air, nitrogen
membawa muatan positif (ion iminium) dan tidak mempunyai hidrogen untuk
dilepaskan; sebagai gantinya, sebuah proton lepas dari karbon di sebelahnya,
sehingga terbentuk ikatan \ce{C=C}.
\end{solution}

\begin{solution}{exo:b2:amines:7}
Benzaldehida dengan propan-2-amina, atau propanon dengan benzilamina,
masing-masing dengan \ce{NaBH3CN} pada pH sekitar 6.
\end{solution}

\begin{solution}{exo:b2:amines:8}
Propiofenon (1-fenilpropan-1-on), \ce{C6H5COCH2CH3}: satu adisi menghasilkan
anion \omterm{def:b2:amines:imine}{imina}, yang dihidrolisis menjadi keton.
\end{solution}

\begin{solution}{exo:b2:amines:9}
Kalium ftalimida + 1-bromoheksana: \textit{N}-heksilftalimida; lalu hidrazina
(atau hidrolisis) membebaskan heksan-1-amina. Nitrogen ftalimida yang
teralkilasi tidak mempunyai pasangan elektron bebas yang tersedia (pasangan
itu terdelokalisasi pada dua karbonil) dan tidak lagi mempunyai \ce{N-H}:
nitrogen itu tidak dapat dialkilasi untuk kedua kalinya.
\end{solution}

\begin{solution}{exo:b2:amines:10}
Air brom membrominasi anilina pada ketiga posisi orto dan para terhadap
gugus amino: 2,4,6-tribromoanilina. \omterm{def:b2:amines:diazonium}{Diazotisasi}, lalu \ce{H3PO2}, menggantikan
gugus amino dengan H: ketiga brom kini saling berposisi 1,3,5.
\end{solution}

\begin{solution}{exo:b2:amines:11}
4-Nitroanilina + \ce{NaNO2} + \ce{HCl} pada 0--\qty{5}{\degreeCelsius}:
4-nitrobenzenadiazonium klorida. Kopling dengan fenol dalam larutan yang
basa lemah (fenoksida, sangat teraktivasi) pada posisi para terhadap
\ce{O-}: 4-[(4-nitrofenil)diazenil]fenol, \ce{O2N-C6H4-N=N-C6H4-OH}, sebuah
zat warna jingga.
\end{solution}

\begin{solution}{exo:b2:amines:12}
Gugus trimetilamonium adalah gugus pergi yang buruk dan meruah, dan membuat
hidrogen-hidrogen $\beta$ pada sisi \ce{CH3} lebih asam dan lebih mudah
dijangkau; basa mengambil hidrogen dari karbon yang kurang tersubstitusi:
but-1-ena, alkena yang kurang tersubstitusi (produk Hofmann).
\end{solution}

\begin{solution}{pb:b2:amines:1}
\textbf{1.} $\mathrm{{}^-O_3S{-}C_6H_4{-}NH_3^+}$: gugus sulfonat yang asam telah memprotonasi amina.
\textbf{2.} Zwitterion itu sukar larut dalam air; karbonat mendeprotonasi
gugus amonium, sehingga terbentuk natrium sulfanilat yang larut.
\textbf{3.} \ce{NaNO2 + HCl -> HNO2 + NaCl}.
\textbf{4.}
\[
  \mathrm{{}^-O_3S{-}C_6H_4{-}NH_2 + HNO_2 + H^+ \to {}^-O_3S{-}C_6H_4{-}N_2^+ + 2\,H_2O} .
\]
\textbf{5.} Garam diazonium terurai (menjadi fenol dan \ce{N2}) jika hangat;
asam nitrit juga terurai.
\textbf{6.} $5.00/173.19 = \qty{28.9}{mmol}$: \qty{28.9}{mmol} \ce{NaNO2}, yaitu \qty{1.99}{g}.
\textbf{7.} Asam nitrit berlebih akan bereaksi dengan pasangan kopling;
kelebihan itu dideteksi dengan kertas kanji--iodida dan dimusnahkan dengan
urea atau asam sulfamat.
\textbf{8.} Cincinnya sangat teraktivasi oleh gugus dimetilamino, sebuah
donor yang kuat.
\textbf{9.} Para terhadap \ce{N(CH3)2}: pengarahan orto/para, dan
posisi-posisi orto terhalang oleh gugus-gugus metil.
\textbf{10.} $\mathrm{{}^-O_3S{-}C_6H_4{-}N{=}N{-}C_6H_4{-}N(CH_3)_2}$.
\textbf{11.} Dalam asam kuat, gugus dimetilamino terprotonasi: \ce{-NH(CH3)2+}
mendeaktifkan cincin dan kopling berhenti.
\textbf{12.} Produknya terbentuk dalam bentuk asamnya yang terprotonasi dan
sebagian larut; dalam basa, natrium sulfonatnya, yang kurang larut dalam
larutan bergaram, mengendap.
\textbf{13.} Muatan positifnya terdelokalisasi ke cincin dan nitrogen
ujungnya merupakan pusat elektrofilik yang buruk; hanya cincin yang
diaktifkan oleh \ce{-O-}, \ce{-OH}, atau \ce{-NR2} yang bereaksi.
\textbf{14.} Kuning jingga dalam basa, merah dalam asam.
\textbf{15.} Pada salah satu nitrogen gugus azo (bentuk terprotonasinya
distabilkan oleh konjugasi dengan gugus dimetilamino).
\textbf{16.} Sekitar $\mathrm pK_a \pm 1$: kira-kira dari 3.1 sampai 4.4, di sekitar $\mathrm pK_a = 3.47$.
\textbf{17.} Dua cincin dan gugus azo membentuk sistem terkonjugasi yang
panjang: celah antara tingkat terisi tertinggi dan tingkat $\pi$ kosong
terendah kecil (\cref{prop:b2:huckel:gap}) dan serapannya terletak di daerah
tampak.
\textbf{18.} Protonasi mengubah sebaran muatan dalam sistem $\pi$ dan makin
mempersempit celah itu: serapannya bergeser ke panjang gelombang yang lebih
panjang, dan warnanya dari kuning menjadi merah.
\textbf{19.} \qty{173.19}{g/mol}.
\textbf{20.} \qty{327.33}{g/mol}.
\textbf{21.} \qty{28.87}{mmol}.
\textbf{22.} $0.02887 \times 327.33 = \qty{9.45}{g}$.
\textbf{23.} Penguraian garam diazonium, kopling yang tidak sempurna,
kelarutan produk dalam larutan induk dan air cucian.
\textbf{24.} $0.80 \times 9.45 = \boldsymbol{\qty{7.56}{g}}$ metil jingga.
\end{solution}
